El plan de validación de usuarios se incluyó como otra validación a la verificación del sistema en la que se busca poder evaluar si Thary es un sistema intuitivo y fácil de usar para los usuarios finales, y que la información sea fácil de comprender e interpretar, de manera que esta le ayude al usuario a tomar decisiones.

La validacccion se realiza para los dos perfiles de usuario implementados: Administrador y Personal de Inventario. 

Para su realización, se va a determinar si el usuario puede completar el flujo completo de uso del sistema sin asistencia adicional más allá de la que brinda la propia interfaz:

\begin{itemize}
    \item \textbf{Administrador}: inicia sesión, carga y procesa el historial de consumo, consulta el panel técnico de comparación de modelos, crea una cuenta de personal de inventario, y revisa la bitacora de actividad de su empresa.
    \item \textbf{Personal de Inventario}: inicia sesión, consulta el dashboard de demanda pronosticada, el estado del inventario, el calendario de pedidos y el presupuesto disponible, confirma un pedido sugerido a partir de las recomendaciones generadas, y revierte dicho pedido para verificar que la acción se pueda deshacer correctamente. Finalmente, consulta su propio perfil.
\end{itemize}

Durante este recorrido se observa si el usuario logra completar cada tarea sin asistencia adicional a la propia interfaz, prestando especial atención a los siguientes aspectos:

\begin{itemize}
    \item Si las alertas de reposición y las recomendaciones de compra son comprensibles sin necesidad de conocimientos técnicos de modelado predictivo.
    \item Si el usuario logra distinguir con claridad las pantallas y funcionalidades disponibles según su rol.
    \item Si el flujo de confirmar y, de ser necesario, deshacer un pedido resulta intuitivo y sin ambiguedad sobre el estado final del pedido.
    \item Si la información del calendario de pedidos y del costo de ruptura (stockout) es interpretada correctamente como una alerta de urgencia, y no como un dato meramente informativo.
\end{itemize}

Al finalizar el recorrido, la percepción del usuario se recopila mediante una breve conversación en la que se le pregunta si encontró alguna pantalla o mensaje confuso y si considera que la información presentada sería útil para tomar decisiones reales de abastecimiento en su contexto laboral. Los resultados de esta validación se documentan en el Capítulo 4.